\documentclass[spanish,openany]{report}
\usepackage{wrapfig}
\input{/home/kuma/git/Clase/template/preamble.tex}
\input{/home/kuma/git/Clase/template/macros.tex}
\input{/home/kuma/git/Clase/template/letterfonts.tex}
\graphicspath{{/home/kuma/git/Clase/Fp_Gs_Robotica/images/}}
\setimagefolder{/home/kuma/git/Clase/Fp_Gs_Robotica/images/}

\date{\today}
% \hbadness=99999
\usepackage{microtype}
% \addbibresource{./bibliografia.bib}
\usepackage{titling}
\usepackage{afterpage}



\begin{document}

%----------------------------------------------------------------------------------------
%	                            PORTADA
%----------------------------------------------------------------------------------------
% \makeportada{}{Daniel Carbajales Soto}{ 1° CFGS Automatización y Robótica Industrial}{}

%----------------------------------------------------------------------------------------
%                               TABLE OF CONTENTS
%----------------------------------------------------------------------------------------
% \maketoc
%----------------------------------------------------------------------------------------
%                                    HEADER 
%----------------------------------------------------------------------------------------
% \header{}{Daniel Carbajales Soto}{1° CFGSARI}
%----------------------------------------------------------------------------------------
%	                            TABLES
%----------------------------------------------------------------------------------------
% this is just a demo in case i need to make a table
%
%
%
% \begin{tabularx}{\linewidth}{|C{2cm}|Y|Y|C{2cm}|}
% \hline
% \makecell{Header 1} & \makecell{Header 2 \\ with linebreak} & \makecell{Header 3} & \makecell{Header 4} \\
% \hline
% Row 1 & This is a long cell text that automatically wraps in the Y column & Another long cell & Short \\
% \hline
% Row 2 & \multirow{2}{*}{Multirow text} & More text & Short \\
% \cline{1-1} \cline{3-4}
% Row 3 & & Extra text & Short \\
% \hline
% \end{tabularx}


%----------------------------------------------------------------------------------------
%	                            DOCUMENT
%----------------------------------------------------------------------------------------
\begin{center}
	\Huge{Ejercicio propuesto pt100}
\end{center}
\vspace{1cm}

\begin{paracol}{2}
\begin{figure}[h]
	\centering
	\includegraphics[width=0.5\textwidth]{./scripts/pt100_plot_kanagawa.pdf}
	\caption*{Curva de características PT100 sin acondicionar}
	\label{fig:fig1}
\end{figure}
\switchcolumn
\begin{center}
		
\begin{tabular}{|c|c|}
\hline
\rowcolor{waveBlue1}
\textcolor{fujiWhite1}{Temperatura (°C) }& \textcolor{fujiWhite1}{Resistencia (\(\Omega\)) }\\
\hline
\rowcolor{waveAqua1}
-100 & 60.26 \\
\hline
\rowcolor{waveAqua1}
0 & 100.00 \\
\hline
\rowcolor{waveAqua1}
100 & 138.51 \\
\hline
\rowcolor{waveAqua1}
200 & 175.86 \\
\hline
\rowcolor{waveAqua1}
300 & 212.05 \\
\hline
\rowcolor{waveAqua1}
400 & 247.09 \\
\hline
\rowcolor{waveAqua1}
500 & 280.98 \\
\hline
\rowcolor{waveAqua1}
600 & 313.71 \\
\hline
\rowcolor{waveAqua1}
700 & 345.28 \\
\hline
\rowcolor{waveAqua1}
800 & 375.70 \\
\hline
\end{tabular}


\vspace{0.75cm}
Valores de la PT100
\end{center}
\end{paracol}
\begin{itemize}
	\item rango:\\
	$-100\si{\degreeCelsius} \text{ -- } 800 \si{\degreeCelsius}$ 
	\item alcance:\\
		\(900\si{\degreeCelsius}\)
	\item sensibilidad:\\
		\(0,34375 \si{\ohm\per\degreeCelsius}\)
	
\end{itemize}





%%%%%%%%%%%%%%%%%%%%%%%%%%%%%%%%%%%%%%%%%%%%%%%%%%%%%%%%%%
%%%%%%%%%%%%%%%%%%%% bibliography %%%%%%%%%%%%%%%%%%%%%%%%
% \nocite{*}                                 %%%%%%%%%%%%%
% \printbibliography                         %%%%%%%%%%%%%
%%%%%%%%%%%%%%%%%%%%%%%%%%%%%%%%%%%%%%%%%%%%%%%%%%%%%%%%%%
%%%%%% Last page in case there is no bibliography %%%%%%%%
% \label{Last Page}                           %%%%%%%%%%%%%
%%%%%%%%%%%%%%%%%%%%%%%%%%%%%%%%%%%%%%%%%%%%%%%%%%%%%%%%%%

\end{document}
